Above that break, therefore, the two-dimensional representation is
trivial on upper wild groups. Its Swan conductor is at most $272$,
and its tame conductor is at most two. Hence
\[
 a_p(r)\le274<300\qquad(p=2,3,5,7).
\]
For $p>7$, $L/\mathbb Q$ is tame because its degree divides $576$,
$A/L$ has good reduction, and $b$ is trivial on wild inertia:
on that pro-$p$ group its square is tame and its remaining image has
order at most two. Thus $a_p(r)\le2$ there; it is zero outside $T$.
Using the coefficient-prime comparison above proves
\[
 N\le210^{300}R^2.
\tag{LR5}
\]

\paragraph{The original constant is preserved.}
Apply the accompanying modular-Jacobian transfer calculation with
$C=210^{300}$ and (LR3)--(LR5). It gives
\[
 h_F(A)<2^{24}210^{3600}R^{24}
 <2^{28824}R^{24}<4608^{165888}R^{24}.
\tag{LR6}
\]
The last comparison uses $210<2^8$, $4608>2^{12}$, and
$28824<1990656=12\cdot165888$. Together with the separately proved
CM bound and geometric height comparison, (LR6) yields exactly
\[
 h(s)\le
 \bigl(2(4608^{165888}+6^{100})B_*^{24}+128\bigr)M(s)^{24}.
\]
No contact, coefficient constant, or radical exponent has been changed.
% UNAPPLIED CANDIDATE: global review and published Pan match are pending.
